\begin{annexitem}{annex:p14.strength.proof-transfer}{thm:p14-strength-proof-transfer}{Direction of a contradiction-preserving proof transfer}
\begin{proof}
Suppose a finite \(U\)-proof of falsity existed. Its image under the
specified transformation would be a finite \(V\)-proof of falsity.
Thus absence of the latter excludes the former. This proves the
external consistency implication and fixes its direction: it runs
from consistency of the target to consistency of the source.

For the internal assertion reason in \(\mathsf B\). Assume
\(\neg\exists q\,P_V(q,\bar c_V)\), and for contradiction assume
\(\exists p\,P_U(p,\bar c_U)\). Existential elimination gives a
temporary witness \(p\) with \(P_U(p,\bar c_U)\). The totality
premise gives \(\exists q\,F(p,q)\). Taking its temporary witness
\(q\), the correctness premise yields \(P_V(q,\bar c_V)\).
Existential introduction contradicts the first assumption. Discharge
the two temporary witness scopes, infer
\(\neg\exists p\,P_U(p,\bar c_U)\), and discharge the first
assumption to obtain the displayed implication.

No representation theorem for \(F\), no derivability conditions,
and no particular induction strength of \(\mathsf B\) are used
in this logical step. Those totality/correctness premises must
be proved separately for an intended numerical transformation.
Uniqueness of the graph output is unnecessary for this implication:
totality and contradiction preservation suffice. A genuine function
graph may additionally have provable uniqueness, but it is not
silently added as a prerequisite here.
\end{proof}
\end{annexitem}

\begin{annexitem}{annex:p14.strength.independence-transfer}{thm:p14-strength-independence-transfer}{Both relative-consistency directions give independence}
\begin{proof}
Assume external consistency of \(T\). The two hypotheses make both
\(T+\varphi\) and \(T+\neg\varphi\) consistent.
If \(T\) proved \(\varphi\), that same finite proof would be
available in \(T+\neg\varphi\), since every original axiom is
still an axiom. The added assumption supplies \(\neg\varphi\),
and negation elimination gives a finite contradiction proof in
that extension. This contradicts its consistency. Hence
\(T\notderives\varphi\).

Conversely, a \(T\)-proof of \(\neg\varphi\) would still be
a proof in \(T+\varphi\). Together with the added \(\varphi\)
it would give a contradiction proof in that extension. Its
consistency excludes this case, so \(T\notderives\neg\varphi\).
These two arguments use different consistency implications.
Neither implication by itself supplies the conclusion of the
other argument. They use weakening of the finite proof's axiom
set and the ordinary contradiction rule; they do not choose an
intended model or decide which sign is true in such a model.
\end{proof}
\end{annexitem}
